\section*{Chapter \ref{ch:g10:chemical-species} --- Chemical Species, Natural and Synthetic}

\begin{solution}{exo:g10:chemical-species:1}
Natural: vanillin from a pod, caffeine from coffee beans. Synthetic:
aspirin, factory vanillin.
\end{solution}

\begin{solution}{exo:g10:chemical-species:2}
They are the same \omterm{def:g6:pure-substances-and-mixtures:species}{chemical species}, and a species has the same
properties whatever its origin.
\end{solution}

\begin{solution}{exo:g10:chemical-species:3}
$R_f = 2.4/6.0 = 0.40$.
\end{solution}

\begin{solution}{exo:g10:chemical-species:4}
It cools the vapour and turns it back into liquid. Entering at the lower
end, the water fills the whole jacket (and meets the hottest vapour last,
cooling best).
\end{solution}

\begin{solution}{exo:g10:chemical-species:5}
\omterm{def:g2:mixing-and-dissolving:dissolve}{Dissolve} the species much better than water; be \omterm{def:g6:solutions-and-solubility:miscible}{immiscible} with water;
be as little hazardous as possible and easy to remove.
\end{solution}

\begin{solution}{exo:g10:chemical-species:6}
Ethanol is \omterm{def:g6:solutions-and-solubility:miscible}{miscible} with water: no layers form, nothing can be separated.
Cyclohexane is \omterm{def:g6:solutions-and-solubility:miscible}{immiscible} with water and \omterm{def:g2:mixing-and-dissolving:dissolve}{dissolves} linalool well.
\end{solution}

\begin{solution}{exo:g10:chemical-species:7}
Water is on top: \qty{1}{mL} of dichloromethane weighs \qty{1.33}{g},
more than \qty{1}{mL} of water, so it sinks.
\end{solution}

\begin{solution}{exo:g10:chemical-species:8}
It melts too low (aspirin: \qty{135}{\celsius}) and over a range: the
sample is impure, or not aspirin.
\end{solution}

\begin{solution}{exo:g10:chemical-species:9}
Linalool and linalyl ethanoate (its two spots are at the heights of the
two references). Linalyl ethanoate climbed higher.
\end{solution}

\begin{solution}{exo:g10:chemical-species:10}
\omterm{def:g10:chemical-species:distillation}{Hydrodistillation}; adding cyclohexane to the distillate and shaking;
separating the layers; \omterm{def:g3:separating-mixtures:evaporate}{evaporating} the cyclohexane.
\end{solution}

\begin{solution}{exo:g10:chemical-species:11}
Pencil does not \omterm{def:g2:mixing-and-dissolving:dissolve}{dissolve} in the \omterm{def:g10:chemical-species:tlc}{eluent}; above the \omterm{def:g10:chemical-species:tlc}{eluent}, the spots are
carried up the plate instead of dissolving into the tank.
\end{solution}

\begin{solution}{exo:g10:chemical-species:12}
Cyclohexane boils at \qty{81}{\celsius}, linalool at \qty{198}{\celsius}:
gentle heating drives off the cyclohexane while the linalool stays.
\end{solution}

\begin{solution}{exo:g10:chemical-species:13}
The sample is probably that species (same $R_f$ as the natural one on
the same plate) and probably pure (one spot, sharp melting point).
\end{solution}

\begin{solution}{exo:g10:chemical-species:14}
$0.40 \times 5.5 = \qty{2.2}{cm}$. The other at $0.75 \times 5.5 =
\qty{4.1}{cm}$ (to the nearest mm), so about \qty{1.9}{cm} above A.
\end{solution}

\begin{solution}{exo:g10:chemical-species:15}
\qty{200}{mL} can \omterm{def:g2:mixing-and-dissolving:dissolve}{dissolve} $1.6 \times 0.200 = \qty{0.32}{g}$: with
\qty{0.20}{g} it is not saturated; it could \omterm{def:g2:mixing-and-dissolving:dissolve}{dissolve} \qty{0.12}{g} more.
The \qty{0.20}{g} dissolved would be lost with the water; cyclohexane,
which \omterm{def:g2:mixing-and-dissolving:dissolve}{dissolves} linalool much better, takes most of it back.
\end{solution}

\begin{solution}{pb:g10:chemical-species:1}
\textbf{1.} Its steam carries the fragrant species out of the flowers and
into the condenser.

\textbf{2.} It cools the vapour back into a liquid; the cold water enters
at the lower end.

\textbf{3.} \omterm{def:g6:solutions-and-solubility:miscible}{Immiscible}; less dense than water (it floats).

\textbf{4.} The layer is very thin and part of the oil is dissolved in the
water or dispersed in it: an \omterm{def:g10:chemical-species:extraction}{extraction} recovers it.

\textbf{5.} It \omterm{def:g2:mixing-and-dissolving:dissolve}{dissolves} linalool well, is \omterm{def:g6:solutions-and-solubility:miscible}{immiscible} with water, and
boils at a low temperature (\qty{81}{\celsius}), so it is easy to remove.

\textbf{6.} No: ethanol \omterm{def:g2:mixing-and-dissolving:mix}{mixes} with water and no layers form.

\textbf{7.} On top (\qty{0.78}{g} per mL, less than water). With
dichloromethane (\qty{1.33}{g} per mL) the \omterm{def:g6:solutions-and-solubility:solute}{solvent} layer would be at the
bottom.

\textbf{8.} No flame nearby (flammable); work under a fume hood with
gloves, and collect the waste instead of pouring it down the sink
(toxic to aquatic life).

\textbf{9.} Cyclohexane boils at \qty{81}{\celsius}, linalool at
\qty{198}{\celsius}: gentle heating removes the \omterm{def:g6:solutions-and-solubility:solute}{solvent} only.

\textbf{10.} $R_f(\text{L}) = 2.4/6.0 = 0.40$; $R_f(\text{A}) = 4.5/6.0
= 0.75$.

\textbf{11.} Linalool and linalyl ethanoate.

\textbf{12.} No: it contains at least two species; it is a \omterm{def:g6:pure-substances-and-mixtures:mixture}{mixture}.

\textbf{13.} So that all spots are developed in the same conditions:
only then can heights be compared.

\textbf{14.} Its $R_f$ equals that of linalool in the same conditions: it
is probably linalool, identical to the \omterm{def:g10:chemical-species:natural}{natural species}.

\textbf{15.} $1.0 \times 0.90 = \qty{0.90}{g}$.

\textbf{16.} $0.90/100 = \qty{0.90}{\percent}$.

\textbf{17.} $100/0.009 \approx 11\,000$\,g, about \qty{11}{kg} of
flowers.

\textbf{18.} It is extracted from a plant, without \omterm{def:g5:irreversible-changes:chemical-change}{chemical change}. But its
linalool is the same species as synthetic linalool: same formula, same
properties.

\textbf{19.} \qty{0.90}{g} of essential oil from \qty{100}{g} of flowers.
\end{solution}
